%% kapa: exam=1 round=28 year=2017 period=2 form=A subject=appraisal_law subject_ko=감정평가관계법규 seq=20 number=20
%% source: Q-Net 감정평가사 기출문제 (cst00309) artlSeq=5199603 file=extracted/1008/5199603/2017년 제28회 감정평가사 1차시험 문제지/2교시 A형.pdf
%% conversion: kordoc 4.24.51 -> tools/gen1.py (auto); review=proofread-sonnet; flags=none
\begin{question}{20}
국유재산법령상 행정재산에 관한 설명으로 옳지 않은 것은?
\begin{choices}
  \item 중앙관서의 장은 행정재산을 효율적으로 관리하기 위하여 필요하면 국가기관 외의 자에게 그 재산의 관리를 위탁할 수 있다.
  \item 행정재산을 사용허가하려는 경우 수의(隨意)의 방법으로는 사용허가를 받을 자를 결정할 수 없다.
  \item 행정재산의 관리위탁을 받은 자는 미리 해당 중앙관서의 장의 승인을 받아 위탁 받은 재산의 일부를 다른 사람에게 사용ㆍ수익하게 할 수 있다.
  \item 중앙관서의 장은 건물 등을 신축하여 기부채납을 하려는 자가 신축기간에 그 부지를 사용하는 경우 그 사용료를 면제할 수 있다.
  \item 중앙관서의 장은 행정재산의 사용허가를 철회하려는 경우에는 청문을 하여야 한다.
\end{choices}
\end{question}
\answer{2}
